\documentclass[11pt,a4paper]{article}
\input{tutorial_style}

\begin{document}
\tutorialtitle{Code Running Guide}{How to Run and Extend the Note 1 Python Examples}{Installation, smoke tests, training diagnostics, outputs, and troubleshooting}{Companion to the cyber control and game-learning guide}
\tableofcontents
\newpage
\lhead{\small Note 1 Code Running Guide}

\section{What This Guide Covers}
This guide is a standalone manual for the Python examples. It explains how to set up the environment, run smoke tests, run longer examples, interpret outputs, and extend the code to more complex cyber models.

\section{Folder Structure}
The repository is organized around one readable path from model to experiment:
\begin{verbatim}
.
  requirements.txt
  README.md
  README.md
  scripts/
    run_smoke_tests.sh
    generate_figures.py
    run_training_iterations.py
  src/
    cyber_dynamics.py
    sampled_continuous_impulse_env.py
    fbsm_malware_baseline.py
    ddqn_cyber_defense.py
    madrl_ctde_parameterized_game.py
    node_level_robustness.py
    evaluation_metrics.py
  artifacts/experiments/
    OUTPUT_PREVIEW.md
    training_summary.md
    policy_comparison_metrics.csv
    game_response_metrics.csv
    node_level_robustness_metrics.csv
  docs/assets/
    action_timing.png
    node_level_learning_advantage.png
\end{verbatim}

\section{Installation}
Create a Python virtual environment and install dependencies:
\begin{lstlisting}[style=pythonstyle]
python -m venv .venv
source .venv/bin/activate          # Windows: .venv\Scripts\activate
pip install --upgrade pip
pip install -r requirements.txt
\end{lstlisting}
The examples require only NumPy, PyTorch, and optionally Matplotlib. CUDA is optional.

\section{Quick Smoke Test}
Run all scripts in quick mode:
\begin{lstlisting}[style=pythonstyle]
bash scripts/run_smoke_tests.sh
\end{lstlisting}
A smoke test checks syntax and basic data flow. It does not train high-quality policies or prove game-theoretic convergence.

\section{Normal Runs}
\subsection{Classical FBSM baseline}
\begin{lstlisting}[style=pythonstyle]
python src/fbsm_malware_baseline.py --max-iter 100
\end{lstlisting}
Expected output: iteration logs, peak infection, objective value, final state, and mean control.

\subsection{DDQN defender}
\begin{lstlisting}[style=pythonstyle]
python src/ddqn_cyber_defense.py --episodes 300
\end{lstlisting}
Expected output: training return, validation return, and exploration rate. Improvement is stochastic; run multiple seeds.

\subsection{MADRL attacker-defender game}
\begin{lstlisting}[style=pythonstyle]
python src/madrl_ctde_parameterized_game.py --episodes 200
\end{lstlisting}
Expected output: defender and attacker episode returns. This file is a compact CTDE teaching example. For research, extend it with MAPPO clipping, GAE, opponent pools, and cross-play evaluation.

\subsection{Static figure generation}
\begin{lstlisting}[style=pythonstyle]
python scripts/generate_figures.py
\end{lstlisting}
Expected output: refreshed PNG figures under \texttt{docs/assets/}, including neural architecture, action timing, FBSM, sampled-flow rollout, and policy-comparison plots.

\subsection{Longer training diagnostics}
\begin{lstlisting}[style=pythonstyle]
python scripts/run_training_iterations.py --episodes 180
\end{lstlisting}
Expected output:
\begin{itemize}
\item CSV histories under \texttt{artifacts/experiments/};
\item \texttt{OUTPUT\_PREVIEW.md} and \texttt{training\_summary.md};
\item the training-diagnostics, game-response, and node-level epidemic-robustness PNG figures.
\end{itemize}

\section{How to Interpret Outputs}
\begin{longtable}{p{0.24\textwidth}p{0.66\textwidth}}
\toprule
Output & Interpretation \\
\midrule
Training return & Cumulative reward during learning. It is noisy and should not be the only metric. \\
Validation return & Performance of the current policy without exploration against a fixed scenario. \\
Peak infected & Maximum compromised population; useful cyber-risk metric. \\
Area under $I(t)$ & Total compromise burden over time; often more informative than final infection. \\
Defender cost & Negative cumulative defender reward; lower is better when exposure is also controlled. \\
Impulse events & How often immediate jump actions, such as isolation, were used. \\
Game response matrix & Defender-policy performance against several attacker strategies. \\
Node-level epidemic robustness plot & Parameter-mismatch stress test: nominal FBSM open-loop schedule versus DDQN aggregate feedback on a stochastic graph where each node has a local S/I/R state. \\
\bottomrule
\end{longtable}

\section{Extending the Code}
\subsection{From simple malware to APT}
Add a deception action and modify the infection rate during the affected decision interval:
\begin{equation}
\beta SI\quad \rightarrow\quad \beta(1-\chi\eta_k)VC,\qquad 0\leq\eta_k\leq 1.
\end{equation}
Here $\eta_k$ is the deception intensity selected at action epoch $t_k$. Add actions such as deceive, isolate, and stealth. The provided environment already contains these mechanisms.

\subsection{From compartment model to degree-based model}
Replace $S,I,R$ by degree-indexed compartments. For example,
\begin{equation}
x=[S_1,I_1,R_1,\ldots,S_K,I_K,R_K].
\end{equation}
The infection pressure should use the degree-weighted infected probability.

\subsection{From scripted attacker to learned attacker}
Start with \texttt{ddqn\_cyber\_defense.py}. Replace the scripted attacker with an actor network. Then use \texttt{madrl\_ctde\_parameterized\_game.py} as the training template.

\subsection{From one attacker to a response matrix}
Use \texttt{evaluation\_metrics.py} to add attacker policies, then rerun the training-diagnostics script.  The output \texttt{game\_response\_metrics.csv} is the first step toward cross-play and exploitability analysis.

\subsection{From tabular response to larger games}
For larger games, keep the environment contract explicit: observe at action point $t_k$, apply jumps only when intended, integrate over $[t_k,t_{k+1})$, and report cyber metrics alongside rewards. A fixed interval is convenient, but a nonuniform action schedule is valid if it is recorded and used consistently. Then add opponent pools, multiple random seeds, and unilateral-deviation checks.

\section{Troubleshooting}
\begin{itemize}
\item If DDQN is unstable, reduce reward scale, reduce learning rate, enlarge replay buffer, and update target network less frequently.
\item If MADRL cycles, freeze one agent for several updates, add opponent pools, and evaluate cross-play matrices.
\item If a learned policy looks good in reward but bad in cyber metrics, inspect the reward formula and compare cumulative compromise, peak compromise, and defender cost.
\item If an action label is unclear, rename it in \texttt{evaluation\_metrics.py}; figures and CSVs will then inherit the clearer label.
\item If outputs are missing, run the workflow in order: smoke tests, static figure generation, then longer training diagnostics.
\end{itemize}

\section{Reproducibility Checklist}
For every experiment, record:
\begin{itemize}
\item model equations and parameter values;
\item random seed and initial states;
\item action space and reward formula;
\item ODE solver, action schedule, and solver step size;
\item network architecture and optimizer;
\item number of episodes or iterations;
\item all baselines and attacker strategies used for comparison;
\item output metrics: cumulative compromise, peak compromise, final compromise, defender cost, impulse count, and game-response table.
\end{itemize}

\end{document}
